\subsection{トピックと参考文献の対応}
この表は、SWEBOK が示すトピック対参考資料の行列を、学習用に簡略化したものである。第7章では、Grubb and Takang の Software Maintenance: Concepts and Practice が中心的な参考資料として扱われる。\par
\par\smallskip\noindent\refstepcounter{table}\label{tab:ch07-15}表 \thetable: トピックと参考文献の対応におけるトピック・主な参考資料\par\smallskip
表 \thetable は、トピックと参考文献の対応におけるトピック・主な参考資料を比較・確認しやすい形で整理したものである。\par\smallskip
\begin{longtable}{p{0.30\linewidth}p{0.64\linewidth}}
\toprule
トピック & 主な参考資料 \\
\midrule
\endhead
1. ソフトウェア保守の基礎 & ISO/IEC/IEEE 14764; Grubb and Takang \\
1.1 定義と用語 & ISO/IEC/IEEE 14764; Grubb and Takang \\
1.2 保守の性質 & Grubb and Takang \\
1.3 保守が必要な理由 & Grubb and Takang \\
1.4 保守費用の大きさ & Grubb and Takang; Abran and Nguyenkim \\
1.5 ソフトウェア進化 & Lehman; Grubb and Takang \\
1.6 保守分類 & ISO/IEC/IEEE 14764; Grubb and Takang \\
2.1 技術的課題 & Grubb and Takang; ISO/IEC/IEEE 14764 \\
2.1.1 限られた理解 & Grubb and Takang \\
2.1.2 テスト & Grubb and Takang; Software Testing KA \\
2.1.3 影響分析 & ISO/IEC/IEEE 14764; Grubb and Takang \\
2.1.4 保守性 & ISO/IEC/IEEE 14764; Software Quality KA \\
2.2 管理上の課題 & Grubb and Takang; April and Abran \\
2.3 保守費用 & ISO/IEC/IEEE 14764; Grubb and Takang \\
2.4 保守測定 & Grubb and Takang; Software Engineering Management KA \\
3. 保守プロセス & ISO/IEC/IEEE 14764; ISO/IEC/IEEE 12207; Grubb and Takang \\
4. 保守技法 & Grubb and Takang; Humble and Fairley; Software Quality KA \\
5. 保守ツール & Grubb and Takang; Software Configuration Management KA \\
\bottomrule
\end{longtable}
